
\begin{tikzpicture}[x=1cm,y=1cm,every node/.style={font=\fontsize{10}{13}\selectfont,align=center,inner sep=3pt},line width=.5pt]
\node[anchor=west,font=\fontsize{9.3}{12.1}\selectfont] (chainlabel) at (0,0) {行为结构：};
\node[draw,text width=1.6cm,font=\fontsize{9.3}{12.1}\selectfont] (c0) at (2.9000000000000004,0) {欺骗行为};
\node[draw,text width=2.8cm,font=\fontsize{9.3}{12.1}\selectfont] (c1) at (5.5,0) {对方产生或维持认识错误};
\draw[->] (c0.east)--(c1.west);
\node[draw,text width=2.8cm,font=\fontsize{9.3}{12.1}\selectfont] (c2) at (8.7,0) {对方基于认识错误处分财物};
\draw[->] (c1.east)--(c2.west);
\node[draw,text width=1.8cm,font=\fontsize{9.3}{12.1}\selectfont] (c3) at (11.399999999999999,0) {行为人取得财物};
\draw[->] (c2.east)--(c3.west);
\node[draw,text width=1.8cm,font=\fontsize{9.3}{12.1}\selectfont] (c4) at (13.6,0) {对方遭受财产损失};
\draw[->] (c3.east)--(c4.west);
\node[font=\fontsize{9.3}{12.1}\selectfont] (role0) at (2.9000000000000004,-1.0) {行为人};
\node[font=\fontsize{9.3}{12.1}\selectfont] (role1) at (5.5,-1.0) {受骗人};
\node[font=\fontsize{9.3}{12.1}\selectfont] (role2) at (8.7,-1.0) {处分人};
\node[font=\fontsize{9.3}{12.1}\selectfont] (role3) at (11.399999999999999,-1.0) {行为人};
\node[font=\fontsize{9.3}{12.1}\selectfont] (role4) at (13.6,-1.0) {受害人};

\end{tikzpicture}

